\expandafter\def\csname HistoricalNoteText550\endcsname{Gelegentlich (wenn auch äußerst selten) sind uns allerdings Ausnahmefälle begegnet, in denen sich selbst in so kurzen Zwischenzeiten unverkennbar Spuren einer Beschleunigung oder Verlangsamung der Änderung nachweisen ließen. Dieser Gegenstand soll künftig ausführlich behandelt werden,}
\expandafter\def\csname HistoricalNoteText559\endcsname{Zu dieser Einrichtung sei hier noch einiges bemerkt. Eine Marke für die erwähnte Prüfung ist als unerlässliche Voraussetzung für die Zuverlässigkeit der Beobachtungen anzusehen. Beim Bau eines neuen Gebäudes muss man sie daher gebührend berücksichtigen. Vor dem Bau unseres magnetischen Observatoriums war auch erwogen worden, ob man diese Marke besser auf einem eigenen Sockel mit gesondertem Fundament im Inneren des Saals anbringen sollte als an der Wand, die außen der Witterung ausgesetzt ist. Man entschied sich für die Wand, da man sonst entweder den Abstand des Beobachters von der Nadel hätte verringern, den Vorteil, Marke und Skala bei derselben Okularstellung deutlich zu sehen, hätte aufgeben oder den Saal noch länger hätte machen müssen. was an dem vorgesehenen Ort nicht einmal möglich gewesen wäre. Ein künstlicher Ersatz für eine Marke wurde aus mehreren Gründen verworfen. Auch die Sorge, dass sich der absolute Ort der Marke durch die Einwirkung der Witterung auf die Wand merklich ändern könnte, hielt man bei solider Bauausführung und der sehr geringen Höhe der Marke über der Grundmauer für wenig erheblich, zumal man die Winkelmessung zwischen der Marke und dem durch das Nordfenster sichtbaren Kirchturm jederzeit wiederholen kann. Eine dreijährige Erfahrung spricht jetzt für die Richtigkeit dieser Ansicht.}
